\chapter{The Full Multiscale Graph VFE}
\label{ch:full-graph-meta-agent-vfe}

The exact variational free energy was already defined in \Cref{ch:elbo},
extended to one normalized interacting population in
\Cref{ch:local-collective-elbo}, and transported through one declared coarse
channel in \Cref{ch:agent-network-rg}. This chapter does not replace those
results. It asks what additional probabilistic structure is required to put
an adaptive two-channel graph, emergent partitions, and parent-to-child
influence inside one multiscale VFE.

The answer has two levels. The exact answer is always the VFE of one
normalized generative law and one normalized recognition law. A sum of self,
edge, partition, and cross-scale terms is exact only when it is obtained by
factorizing those two laws. Without such a factorization the same sum can be a
useful composite descent potential, but it is not automatically an ELBO and
does not inherit an evidence interpretation. \status{ESTABLISHED}

\section{The full answer before decomposition}
\label{sec:full-graph-vfe-global}

Fix structural data $X$, an admitted observation
$o\in\mathsf O_{\theta,X}^{\mathrm{reg}}$, and a finite scale depth $S$. Let
$W$ contain every random object retained by the multiscale model: all belief
and model coordinates at scales $0,\ldots,S$, edge and connection variables,
partition or membership variables, retained holonomy marks, and auxiliary
latents. Declare one normalized generative kernel
$\mathbb P_\theta(do,dW\given X)$ before recognition and one normalized
recognition kernel $\mathbb Q_\phi(dW\given o,X)$. With selected posterior
$\boldsymbol\Pi_{\theta,o,X}$ and evidence density
$p_\theta(o\given X)\in(0,\infty)$, the full multiscale VFE is
\begin{equation}
\boxed{
\mathcal F_{\mathrm{tower}}^{\mathrm{ext}}
[\mathbb Q_\phi;X,o]
=-\log p_\theta(o\given X)
+\KL\left(
\mathbb Q_\phi(\cdot\given o,X)
\middle\Vert
\boldsymbol\Pi_{\theta,o,X}
\right).}
\label{eq:full-graph-vfe-global}
\end{equation}
Equivalently, on the finite-density domain,
\begin{equation}
\mathcal F_{\mathrm{tower}}
=\E_{\mathbb Q_\phi}
\left[
\log\frac{q_\phi(W\given o,X)}{p_\theta(o,W\given X)}
\right].
\label{eq:full-graph-vfe-density}
\end{equation}
Equations~\eqref{eq:full-graph-vfe-global} and
\eqref{eq:full-graph-vfe-density} are the complete VFE. Every later term is an
accounting decomposition of this object or is labeled as a proposed composite
potential. \status{ESTABLISHED}

\section{A normalized hierarchical graph model}
\label{sec:full-graph-vfe-joint}

At scale $s$, let $V_s$ be the vertex set and let
$Z_s=(Z_{i,s})_{i\in V_s}$ contain the full probabilistic coordinates of that
scale, including belief and model presentations rather than only their
marginals. Let $G_s$ contain the directed graph data
$(\beta^s,\gamma^s,\Omega^{b,s},\Omega^{m,s})$, and let $R_s$ be a possibly
stochastic membership relation from $V_s$ to $V_{s+1}$. The structural
configuration $X$ and observation $O$ stay outside the coarse channels.

A normalized top-down generative factorization is
\begin{align}
\mathbb P_\theta(do,dW\given X)
={}&L_\theta(do\given Z_0,X)P_S(dZ_S,dG_S\given X)
\nonumber\\
&\times\prod_{s=0}^{S-1}
P_R^s(dR_s\given Z_{s+1},G_{s+1},X)
P_G^s(dG_s\given Z_{s+1},R_s,X)
\nonumber\\
&\times\prod_{s=0}^{S-1}
K_\downarrow^s(dZ_s\given Z_{s+1},R_s,G_s,X).
\label{eq:full-graph-generative-factorization}
\end{align}
Every factor is a normalized measurable kernel. In particular,
$K_\downarrow^s$ is the mathematical location of downward influence: a parent
belief or model changes the conditional prior of its children. The bottom-up
coarse map belongs to recognition or to an induced pushforward; it is not
inserted as a second incompatible generative arrow at the same time slice.
\status{HYPOTHESIS}

For an exact chain-rule decomposition, disintegrate recognition along the same
ordering:
\begin{align}
\mathbb Q_\phi(dW\given o,X)
={}&Q_S(dZ_S,dG_S\given o,X)
\nonumber\\
&\times\prod_{s=0}^{S-1}
Q_R^s(dR_s\given Z_{s+1},G_{s+1},o,X)
Q_G^s(dG_s\given Z_{s+1},R_s,o,X)
\nonumber\\
&\times\prod_{s=0}^{S-1}
Q_s(dZ_s\given Z_{s+1},R_s,G_s,o,X).
\label{eq:full-graph-recognition-factorization}
\end{align}
This is a disintegration of a correlated law, not a mean-field assumption. A
bottom-up recognition network may parameterize the same joint; the written
ordering is only an accounting choice. \status{HYPOTHESIS}

\theoremheading{Exact hierarchical decomposition}{thm:full-graph-vfe-decomposition}
Assume the normalized factorizations above, common domination, and the
integrability needed for the finite log-density form. Then
\begin{align}
\mathcal F_{\mathrm{tower}}
={}&-\E_{\mathbb Q_\phi}\log L_\theta(o\given Z_0,X)
+\KL(Q_S\Vert P_S)
\nonumber\\
&+\sum_{s=0}^{S-1}\E_{\mathbb Q_\phi}
\KL\left(Q_R^s(\cdot\given Z_{s+1},G_{s+1},o,X)
\middle\Vert P_R^s(\cdot\given Z_{s+1},G_{s+1},X)\right)
\nonumber\\
&+\sum_{s=0}^{S-1}\E_{\mathbb Q_\phi}
\KL\left(Q_G^s(\cdot\given Z_{s+1},R_s,o,X)
\middle\Vert P_G^s(\cdot\given Z_{s+1},R_s,X)\right)
\nonumber\\
&+\sum_{s=0}^{S-1}\E_{\mathbb Q_\phi}
\KL\left(Q_s(\cdot\given Z_{s+1},R_s,G_s,o,X)
\middle\Vert K_\downarrow^s(\cdot\given Z_{s+1},R_s,G_s,X)\right).
\label{eq:full-graph-vfe-factorized}
\end{align}
Each conditional divergence is averaged over its displayed parent variables
under $\mathbb Q_\phi$. These are the observation, top-prior, partition,
graph, and parent--child contributions. \status{ESTABLISHED}

\noindent\emph{Proof.}
Substitute the two factorizations into the Radon--Nikodym ratio in
\eqref{eq:full-graph-vfe-density} and apply the conditional relative-entropy
chain rule. The extended-real version must use the measure-level chain rule
without subtracting infinities. \hfill$\square$

A deterministic same-time parent $Z_{s+1}=C_s(Z_s)$ and an independently
sampled parent that generates $Z_s$ are not two free descriptions of one
joint. One may use the hierarchy above, a deterministic pushforward, or a
normalized undirected compatibility factor. Reciprocal same-time influence
requires the last option and its partition function. Delayed Wheelerian
feedback instead lets $Z_{s+1}(t)$ parameterize the child kernel at
$t+\Delta t$; that is a dynamical extension, not a static ELBO identity.
\status{NOT-CLAIMED}

\section{The scale-local two-channel graph energy}
\label{sec:full-graph-vfe-scale-local}

For arbitrary directed, generally non-flat transports, define
\begin{align}
D_{ij}^{b,s}
&:=\KL\left(q_{i,s}^b\middle\Vert
(\Omega_{ij}^{b,s})_\#q_{j,s}^b\right),
&
D_{ij}^{m,s}
&:=\KL\left(q_{i,s}^m\middle\Vert
(\Omega_{ij}^{m,s})_\#q_{j,s}^m\right).
\label{eq:full-graph-transported-divergences}
\end{align}
The rows $\beta_i^s$ and $\gamma_i^s$ are normalized conditional source laws.
With external receiver occupancies $\alpha_{i,s}^b$ and
$\alpha_{i,s}^m$, define
\begin{equation}
\eta_{ij}^{b,s}=\alpha_{i,s}^b\beta_{ij}^s,
\qquad
\eta_{ij}^{m,s}=\alpha_{i,s}^m\gamma_{ij}^s.
\label{eq:full-graph-edge-event-laws}
\end{equation}
The correct energy-unit row free energies are
\begin{align}
\Phi_{i,s}^b
&=\sum_j\beta_{ij}^sD_{ij}^{b,s}
+\tau_{i,s}^b\KL(\beta_i^s\Vert\pi_i^{b,s}),
\nonumber\\
\Phi_{i,s}^m
&=\sum_j\gamma_{ij}^sD_{ij}^{m,s}
+\tau_{i,s}^m\KL(\gamma_i^s\Vert\pi_i^{m,s}).
\label{eq:full-graph-row-free-energies}
\end{align}
Their conditional minimizers are
\begin{equation}
\beta_{ij}^{s,*}
=\frac{\pi_{ij}^{b,s}e^{-D_{ij}^{b,s}/\tau_{i,s}^b}}
{\sum_k\pi_{ik}^{b,s}e^{-D_{ik}^{b,s}/\tau_{i,s}^b}},
\qquad
\gamma_{ij}^{s,*}
=\frac{\pi_{ij}^{m,s}e^{-D_{ij}^{m,s}/\tau_{i,s}^m}}
{\sum_k\pi_{ik}^{m,s}e^{-D_{ik}^{m,s}/\tau_{i,s}^m}}.
\label{eq:full-graph-gibbs-rows}
\end{equation}
Thus $\beta_{ij}$ and $\gamma_{ij}$ coevolve with the transported divergences.
Omitting the row relative entropy changes the objective. \status{ESTABLISHED}

The quantity $1/\beta_{ij}$ is a useful monotone interaction length only in a
fixed row. It is directed, depends on all competing sources through row
normalization, diverges at zero, and need not obey a triangle inequality. At
the Gibbs optimum the source-relative surprisal
\begin{equation}
\ell_{ij}^{b,s}:=-\tau_{i,s}^b
\log\frac{\beta_{ij}^{s,*}}{\pi_{ij}^{b,s}}
=D_{ij}^{b,s}+\tau_{i,s}^b\log Z_i^{b,s}
\label{eq:full-graph-attention-length}
\end{equation}
is the transported mismatch plus a row-dependent constant. Raw conductance,
a symmetrized resistance, or diffusion time is required for a genuine graph
metric. \status{ESTABLISHED}

If a scale-local prior is a normalized Gibbs law containing the self terms and
$\sum_{ij}\eta_{ij}^{b,s}D_{ij}^{b,s}
+\sum_{ij}\eta_{ij}^{m,s}D_{ij}^{m,s}$, those terms and the global normalizer
belong to the exact VFE. If they are merely appended to an already fixed
state-level VFE, define instead
\begin{align}
\mathcal G_{\mathrm{composite}}
={}&\mathcal F_{\mathrm{population}}
+\sum_s\sum_i\left(
\alpha_{i,s}^b\Phi_{i,s}^b
+\alpha_{i,s}^m\Phi_{i,s}^m\right)
\nonumber\\
&+\sum_s\mathcal F_{\mathrm{partition}}^s
+\sum_s\mathcal F_{\mathrm{cross}}^s
+\sum_s\mathcal R_{\mathrm{connection}}^s.
\label{eq:full-graph-composite-potential}
\end{align}
This is a composite multiscale free-energy potential, not automatically an
evidence bound. \status{DEFINITION}

\section{Non-flat edges and what the parent retains}
\label{sec:full-graph-vfe-nonflat}

The edge transports are not assumed flat. For an oriented cycle
$C=e_1\cdots e_r$, let
\begin{equation}
H_C^x=\Omega_{e_r}^{x}\cdots\Omega_{e_1}^{x},
\qquad x\in\{b,m\}.
\label{eq:full-graph-cycle-holonomy}
\end{equation}
Nonidentity $H_C^x$ is allowed. A holonomy-blind parent law $Q_I^x$ can be
path independent only on the stabilized sector
\begin{equation}
(H_C^x)_\#Q_I^x=Q_I^x
\quad\text{for every retained loop }C.
\label{eq:full-graph-holonomy-stabilization}
\end{equation}
This is not a flatness condition and it does not select $I$. If it fails, the
exact parent must retain a root, path, holonomy representation, or boundary
mark. \status{ESTABLISHED}

For blocks $I,J$, dress each microscopic edge into parent coordinates,
\begin{equation}
\Theta_{ij}^{IJ,x}
=\Omega_{Ii}^{x}\Omega_{ij}^{x}\Omega_{jJ}^{x}.
\label{eq:full-graph-dressed-edge}
\end{equation}
On a non-flat graph the family
$\{\Theta_{ij}^{IJ,x}:i\in I,j\in J\}$ need not collapse to one group
element. Exact coarse data retains its conditional distribution or a mean
together with a declared residual. Penalizing that residual to zero is an
extra modeling preference, not a consequence of gauge covariance.
\status{ESTABLISHED}

\section{Partition variables and graph renormalization}
\label{sec:full-graph-vfe-partitions}

The weights $\beta$ and $\gamma$ do not themselves create a parent random
variable. A hierarchy intended to emerge under the same variational principle
must put $R_s$ inside $\mathbb P$ and $\mathbb Q$. A common finite-temperature
form is
\begin{equation}
\mathcal F_{\mathrm{partition}}^s
=\E_{Q_R^s}[E_{\mathrm{block}}^s(R_s;Z_s,G_s)]
+\tau_R^s\KL(Q_R^s\Vert\Pi_R^s),
\label{eq:full-graph-partition-free-energy}
\end{equation}
provided the normalized Gibbs factor represented by this expression is part
of the generative law. The block energy may trade internal transported
mismatch against boundary retention, evaluator closure, and residual cost.
No canonical selector follows from the global VFE alone. \status{HYPOTHESIS}

Coarse attention is obtained by pushing the joint edge-event law, not by
averaging normalized rows. For deterministic blocks,
\begin{equation}
\eta_{IJ}^{x,s+1}
=\sum_{i\in I}\sum_{j\in J}\eta_{ij}^{x,s},
\qquad
\alpha_I^{x,s+1}=\sum_J\eta_{IJ}^{x,s+1},
\qquad
\beta_{IJ}^{x,s+1}
=\frac{\eta_{IJ}^{x,s+1}}{\alpha_I^{x,s+1}},
\label{eq:full-graph-edge-pushforward}
\end{equation}
on positive parent occupancy. The model channel uses the same construction
with $\gamma$. Stochastic or overlapping membership replaces indicator sums
by expectations under $Q_R^s$. Nested application requires a composition
theorem for membership kernels and retained marks. \status{ESTABLISHED}

\section{How a parent influences its children}
\label{sec:full-graph-vfe-downward}

The normalized kernel $K_\downarrow^s$ supplies the cleanest static answer.
For a block $I$ with children $i\in I$, a factorization such as
\begin{equation}
K_\downarrow^s(dZ_s\given Z_{s+1},R_s,G_s,X)
=\prod_{I\in V_{s+1}}K_I^s
\left(dZ_I^s\given Z_{I,s+1},R_s,G_s,X\right)
\label{eq:full-graph-downward-kernel}
\end{equation}
makes the parent belief and model presentations arguments of the children's
conditional generative prior. The corresponding term is
\begin{equation}
\mathcal F_{\mathrm{cross}}^s
=\E_{\mathbb Q_\phi}
\KL\left(
Q_s(dZ_s\given Z_{s+1},R_s,G_s,o,X)
\middle\Vert
K_\downarrow^s(dZ_s\given Z_{s+1},R_s,G_s,X)
\right).
\label{eq:full-graph-cross-scale-kl}
\end{equation}
Variation with respect to a parent coordinate changes the parent update;
variation with respect to a child coordinate changes the child update. The
optimization is bidirectionally coupled even though the normalized generative
arrow is top-down. This is the precise static sense in which a meta-agent can
influence its sub-agents without allowing a generative factor to read the
posterior. \status{ESTABLISHED}

Calling the optimization parameter physical time requires more. Delayed
feedback may use the inferred parent at time $t$ to modify
$K_\downarrow^s$ at $t+\Delta t$. A same-time reciprocal factor may instead
be placed in a normalized Gibbs joint. Either construction needs work or flux
accounting before it supports Wheelerian nonequilibrium. Static VFE descent
alone proves neither sustained nonequilibrium nor downward causation in
physical time. \status{OPEN}

\section{Bayesian and network renormalization}
\label{sec:full-graph-vfe-literature}

Bayesian renormalization uses Fisher distinguishability on model-parameter
space to define an information scale and separate stiff from sloppy
directions \citep{BermanKlingerStapleton2023}. It can rank which directions
of a parent model presentation remain empirically distinguishable at a
declared observational precision. It does not by itself select graph blocks,
renormalize non-flat transports, or prove closure of the two-channel
interaction law. Its parameter-space orientation must not be identified
silently with sample-space graph diffusion.

Network renormalization addresses the absence of a lattice, momentum, and
geometric length in a heterogeneous graph \citep{gabrielli2025network}.
Laplacian RG offers a diffusion scale and spectral diagnostics
\citep{villegas2023laplacian}; the multiscale model gives a geometry-free
closure family based on additive hidden variables
\citep{Garuccio2023Multiscale}. For the present theory these are candidate
partition and closure mechanisms. They do not automatically handle directed
row-normalized $\beta,\gamma$, two gauge representations, non-flat holonomy,
or simultaneous renormalization of beliefs, models, topology, and dynamics.
The network-RG review identifies the last problem as open.

The proposed synthesis is hybrid: use edge-event laws or raw conductances to
construct a gauge-invariant multiplex graph operator; use spectral or
Bayesian information diagnostics to propose scales and relevant directions;
infer $R_s$ through a normalized partition factor; push the full
probabilistic law and edge-event measures; retain non-flat holonomy or prove
stabilization; and measure failure of the parent family by an explicit closure
residual. \status{HYPOTHESIS}

\section{The next theorem}
\label{sec:full-graph-vfe-target}

Starting at one shared $c_*\in\mathcal C$ with a finite directed multiplex
graph $(\beta_{ij},\gamma_{ij},\Omega_{ij}^b,\Omega_{ij}^m)$, construct a
normalized partition law and parent--child kernels such that descent of
\eqref{eq:full-graph-vfe-global} produces persistent nested blocks, the
pushforward rule \eqref{eq:full-graph-edge-pushforward} composes across scales,
the parent family has a controlled closure residual, and non-flat holonomy is
retained or stabilized without imposing flatness. Then determine whether the
coupled variational dynamics has a separated slow manifold on which
scale-$s+1$ variables feed back consistently into scale-$s$ kernels.

No current result proves that arbitrary descent from arbitrary
$\beta,\gamma$ produces such a hierarchy. The exact VFE tells us how a
proposed hierarchy must be normalized and accounted for; it does not supply
the partition selector, persistence theorem, or dynamical cross-scale
closure. Those are the next proof obligations. \status{OPEN}
